\documentclass[12pt,a4paper]{report}
\usepackage[utf8]{inputenc}
\usepackage[T1]{fontenc}
\usepackage[french]{babel}
\usepackage{lmodern}
\usepackage[a4paper,margin=2.35cm,headheight=18pt]{geometry}
\usepackage{graphicx}
\usepackage{xcolor}
\usepackage{tabularx}
\usepackage{longtable}
\usepackage{booktabs}
\usepackage{array}
\usepackage{enumitem}
\usepackage{fancyhdr}
\usepackage{hyperref}
\usepackage{caption}
\usepackage{float}
\usepackage{pdflscape}
\usepackage[most]{tcolorbox}
\usepackage{tikz}

\definecolor{SmartNavy}{HTML}{173F66}
\definecolor{SmartTeal}{HTML}{2A837C}
\definecolor{SmartRed}{HTML}{D7261E}
\definecolor{SmartViolet}{HTML}{7D60C8}
\definecolor{SmartOrange}{HTML}{EE8626}
\definecolor{SmartGreen}{HTML}{5D9A50}
\definecolor{SmartLight}{HTML}{F6F8FB}

\hypersetup{colorlinks=true, linkcolor=SmartNavy, urlcolor=SmartTeal, citecolor=SmartNavy}
\setlength{\parindent}{0pt}
\setlength{\parskip}{0.62em}
\renewcommand{\arraystretch}{1.18}
\captionsetup{font=small,labelfont=bf,textfont=it}

\pagestyle{fancy}
\fancyhf{}
\lhead{\small Rapport de stage}
\rhead{\small Smart WMS}
\cfoot{\thepage}
\renewcommand{\headrulewidth}{0.4pt}

\newtcolorbox{problematique}{
  colback=SmartLight,
  colframe=SmartNavy,
  boxrule=0.8pt,
  arc=2mm,
  title=Problématique directrice,
  fonttitle=\bfseries
}

\newcommand{\studentName}{FERHAN Abdelali}
\newcommand{\schoolName}{ENIAD -- École Nationale de l'Intelligence Artificielle \& du Digital, Berkane}
\newcommand{\internshipCompany}{Smart Automation Technologies, Tanger}
\newcommand{\internshipPeriod}{Août -- octobre 2026}
\newcommand{\reportTitle}{Conception d'une Architecture Intelligente pour la Transformation d'un Warehouse Management System Traditionnel en Smart WMS}
\newcommand{\reportSubtitle}{Modélisation métier, fonctionnelle et applicative}

\graphicspath{{./}}

\begin{document}

\begin{titlepage}
\centering
\vspace*{1.5cm}
{\Large \schoolName \par}
\vspace{0.4cm}
{\large Filière : Génie informatique \par}
\vspace{1.8cm}
{\Huge\bfseries RAPPORT DE STAGE \par}
\vspace{1.2cm}
{\LARGE\bfseries \reportTitle \par}
\vspace{0.4cm}
{\Large \reportSubtitle \par}
\vfill
\begin{tabular}{rl}
Étudiant : & \studentName \\
Organisme d'accueil : & \internshipCompany \\
Encadrante : & Chaimae Kaina \\
Période : & \internshipPeriod \\
\end{tabular}
\vfill
{\large Année universitaire 2026}
\end{titlepage}

\pagenumbering{roman}
\chapter*{Remerciements}
Je tiens à remercier l'équipe de Smart Automation Technologies pour son accueil, son accompagnement et les échanges qui ont permis de cadrer progressivement le sujet du stage. Je remercie particulièrement mon encadrante en entreprise, Chaimae Kaina, pour ses remarques sur la cohérence métier, la lisibilité des modèles et la nécessité de distinguer clairement le socle WMS des capacités avancées.

Je remercie également l'ENIAD pour le cadre pédagogique de cette expérience, ainsi que toutes les personnes ayant contribué, directement ou indirectement, à l'avancement de ce travail.

\tableofcontents
\listoffigures
\listoftables

\clearpage
\pagenumbering{arabic}
\chapter*{Introduction générale}
La gestion d'un entrepôt ne se limite plus à enregistrer des entrées et des sorties de marchandises. Un entrepôt moderne doit connaître l'état du stock, organiser les emplacements, préparer les commandes, coordonner les opérateurs et les équipements, prendre en compte les contraintes propres aux produits et échanger avec plusieurs systèmes de l'entreprise. Dans ce contexte, le Warehouse Management System (WMS) joue un rôle central : il structure les opérations depuis la réception jusqu'à l'expédition et assure la traçabilité des mouvements.

Un WMS traditionnel couvre déjà les fonctions essentielles : réception, contrôle, mise en stock, gestion des emplacements, inventaire, réapprovisionnement, allocation, picking, packing et expédition. La transformation vers un Smart WMS ne consiste donc pas à remplacer ce socle par une intelligence artificielle généralisée. Elle consiste plutôt à renforcer le système par des capacités activables lorsque le besoin, les données et l'infrastructure le justifient : prévision de charge, optimisation de slotting, détection d'anomalies, supervision énergétique, vision industrielle, RFID ou orchestration avec WCS/WES.

Le stage réalisé au sein de Smart Automation Technologies s'inscrit dans cette logique. L'objectif principal est de concevoir une architecture de référence capable de guider la transformation d'un WMS traditionnel vers un Smart WMS. Le travail porte sur la conception métier, fonctionnelle et applicative logique ; il ne vise pas le développement complet d'un produit WMS industriel.

\begin{problematique}
Comment concevoir une architecture Smart WMS suffisamment complète pour couvrir les fonctions essentielles d'un entrepôt, tout en restant générique, modulaire et adaptable aux contraintes produits, aux données disponibles et au niveau d'automatisation de chaque entreprise ?
\end{problematique}

La démarche adoptée suit une logique SCAN -- PLAN -- ACT. La phase SCAN a permis de comprendre les solutions WMS leaders et les limites d'un WMS classique. La phase PLAN a structuré les besoins, les use cases et l'architecture fonctionnelle cible. La phase ACT prépare la validation du modèle, la modélisation UML et une architecture applicative logique pouvant servir de base à une future implémentation.

\chapter{Présentation de l'organisme d'accueil et cadrage du stage}
\section{Présentation générale de Smart Automation Technologies}
Smart Automation Technologies (SAT) est une startup technologique située à Tanger. Elle intervient dans les domaines de la transformation digitale, de l'architecture des systèmes d'information, de la data science et de l'intelligence artificielle. Son positionnement est celui d'un partenaire technologique capable d'accompagner les organisations depuis la compréhension du besoin jusqu'à la conception de solutions adaptées.

Ce contexte est cohérent avec le sujet du stage, car la transformation d'un WMS vers un Smart WMS exige d'abord une compréhension métier solide. Il ne suffit pas d'ajouter des technologies à un processus existant ; il faut identifier les problèmes opérationnels, les données disponibles, les dépendances avec les systèmes externes et les conditions réelles d'activation des capacités avancées.

\section{Contexte et périmètre du stage}
Le stage couvre la période d'août à octobre 2026 et porte sur la conception d'une architecture intelligente pour transformer un WMS traditionnel en Smart WMS. Le périmètre est volontairement centré sur la modélisation métier, fonctionnelle et applicative logique. Les choix définitifs de technologie, de base de données, d'infrastructure cloud ou d'équipement physique restent hors périmètre à ce stade.

\section{Objectifs}
Les objectifs opérationnels du stage sont les suivants :
\begin{itemize}[leftmargin=1.2cm]
  \item réaliser un benchmark fonctionnel des solutions WMS leaders ;
  \item identifier les fonctions communes et les limites d'un WMS traditionnel ;
  \item formaliser les use cases Core, Smart et Optional ;
  \item proposer une architecture fonctionnelle cible modulaire ;
  \item préparer une modélisation UML cohérente ;
  \item traduire la cible fonctionnelle vers une architecture applicative logique.
\end{itemize}

\section{Livrables}
Les livrables attendus sont un benchmark WMS, une cartographie des use cases, une architecture fonctionnelle cible, des modèles UML et une proposition d'architecture applicative. Le projet logiciel local SmartWMS AI Operations Agent est traité comme une maquette d'aide à la décision : il démontre certains principes d'outillage, de règles déterministes et d'agent conversationnel, sans prétendre constituer un WMS complet.

\chapter{État de l'art et benchmark des solutions WMS}
\section{WMS classique et Smart WMS}
Un WMS classique structure les flux inbound, internes et outbound. Il assure la réception, la mise en stock, l'inventaire, le réapprovisionnement, l'allocation, la préparation, le packing et l'expédition. Un Smart WMS conserve ces responsabilités mais ajoute une capacité d'analyse et d'adaptation : il peut prévoir une charge, détecter une anomalie, optimiser une affectation ou recommander une action.

La différence principale n'est donc pas la présence d'un algorithme. Elle réside dans la capacité du système à transformer des données fiables en décisions opérationnelles exécutables.

\section{Solutions étudiées}
Le benchmark a porté sur Manhattan Active Warehouse Management, SAP Extended Warehouse Management et Oracle Warehouse Management Cloud. L'objectif n'était pas de produire un classement commercial, mais d'extraire des principes de conception applicables à une architecture générique.

\begin{table}[H]
\centering
\caption{Synthèse des apports retenus du benchmark}
\begin{tabularx}{\linewidth}{>{\bfseries}p{3.2cm}X X}
\toprule
Solution & Idée directrice & Apport pour le projet \\
\midrule
Manhattan Active WMS & Unification de l'exécution & Relier stocks, tâches, opérateurs, équipements et décisions dans une boucle courte décision -- action. \\
SAP EWM & Richesse du modèle métier & Donner une place centrale aux objets métier : produit, lot, HU, emplacement, ressource, règle et tâche. \\
Oracle WMS Cloud & Orchestration commande -- wave -- allocation -- tâches & Montrer comment transformer une demande client en travail coordonné et mesurable. \\
\bottomrule
\end{tabularx}
\end{table}

\section{Enseignements pour l'architecture cible}
Trois enseignements structurent la suite du rapport. Premièrement, le socle WMS doit rester stable et complet. Deuxièmement, la donnée de référence est un prérequis de l'intelligence : une optimisation n'a pas de sens si les produits, emplacements, lots et règles sont mal représentés. Troisièmement, les capacités Smart doivent être reliées à la gestion des tâches pour produire une action concrète, et non seulement un indicateur dans un tableau de bord.

\chapter{Analyse des besoins et identification des use cases}
\section{Approche SCAN -- PLAN -- ACT}
La démarche d'identification des besoins a consisté à passer d'une observation générale du domaine WMS vers des use cases traçables. La figure~\ref{fig:scan-plan-act} synthétise cette progression.

\begin{figure}[H]
\centering
\includegraphics[width=0.92\linewidth]{figures/chapitre3/demarche_scan_plan_act.pdf}
\caption{Application de la démarche SCAN -- PLAN -- ACT à l'identification et à la validation des besoins.}
\label{fig:scan-plan-act}
\end{figure}

\section{Logique de classification}
La classification Core / Smart / Optional évite de confondre une fonction indispensable avec une capacité avancée ou dépendante d'une infrastructure. Le Core correspond aux capacités nécessaires au fonctionnement du WMS. Le Smart améliore la décision grâce aux données, à l'optimisation ou au machine learning. L'Optional dépend d'un équipement ou d'un système particulier : caméra 3D, RFID, AMR/AGV, WCS/WES ou BMS/EMS.

\begin{figure}[H]
\centering
\includegraphics[width=\linewidth]{figures/chapitre3/classification_activation.pdf}
\caption{Logique Core / Smart / Optional et principe d'activation selon les données et l'infrastructure.}
\label{fig:classification}
\end{figure}

\section{Cartographie des use cases}
Les tableaux suivants conservent la numérotation validée pendant le stage.

\begin{longtable}{p{2.2cm}p{9.2cm}p{3cm}}
\caption{Use cases Réception}\\
\toprule
ID & Use case & Type \\
\midrule
UC-R01 & Planifier rendez-vous et quai & Core \\
UC-R02 & Réceptionner les marchandises & Core \\
UC-R03 & Identifier le produit & Core \\
UC-R04 & Contrôler la quantité & Core \\
UC-R05 & Contrôler la qualité & Core \\
UC-R06 & Gérer une anomalie de réception & Core \\
UC-R07 & Proposer emplacement de stockage & Core \\
UC-R08 & Prévoir la charge de réception & Smart \\
UC-R09 & Compter les cartons par caméra 3D & Optional \\
UC-R10 & Lire l'étiquette par OCR & Optional \\
UC-R11 & Détecter les défauts par vision & Optional \\
UC-R12 & Suivre le conteneur connecté & Optional \\
\bottomrule
\end{longtable}

\begin{longtable}{p{2.2cm}p{9.2cm}p{3cm}}
\caption{Use cases Stockage et Inventaire}\\
\toprule
ID & Use case & Type \\
\midrule
UC-S01 & Mettre en stock / putaway & Core \\
UC-S02 & Gérer les stocks & Core \\
UC-S03 & Gérer les emplacements & Core \\
UC-S04 & Réaliser l'inventaire tournant & Core \\
UC-S05 & Réapprovisionner la zone picking & Core \\
UC-S06 & Détecter la congestion d'une zone & Smart \\
UC-S07 & Surveiller les produits sensibles & Core conditionnel \\
UC-S08 & Gérer lots / DLC / FEFO & Core conditionnel \\
UC-S09 & Choisir l'emplacement optimal / slotting & Smart \\
UC-S10 & Cibler les emplacements à recompter & Smart \\
UC-S11 & Prédire une rupture avant blocage & Smart \\
UC-S12 & Chercher un emplacement alternatif & Smart \\
UC-S13 & Suivre les palettes par RFID & Optional \\
UC-S14 & Déplacer le stock par AMR / AGV & Optional \\
\bottomrule
\end{longtable}

\begin{longtable}{p{2.2cm}p{9.2cm}p{3cm}}
\caption{Use cases Commandes et Expédition}\\
\toprule
ID & Use case & Type \\
\midrule
UC-E01 & Recevoir la commande client & Core \\
UC-E02 & Planifier la commande & Core \\
UC-E03 & Créer la wave / batch & Core \\
UC-E04 & Allouer le stock & Core \\
UC-E05 & Créer tâches de picking & Core \\
UC-E06 & Effectuer le picking & Core \\
UC-E07 & Effectuer le packing & Core \\
UC-E08 & Charger le camion & Core \\
UC-E09 & Expédier la commande & Core \\
UC-E10 & Guider le picking RF / Pick-to-Light & Optional \\
UC-E11 & Optimiser le parcours picking & Smart \\
UC-E12 & Contrôler le picking scan / poids / vision & Optional \\
UC-E13 & Recommander l'emballage & Smart \\
UC-E14 & Contrôler le chargement scan / caméra & Optional \\
UC-E15 & Détecter commandes à risque de retard & Smart \\
\bottomrule
\end{longtable}

\begin{table}[H]
\centering
\caption{Synthèse des catégories et critères d'activation par domaine}
\begin{tabularx}{\linewidth}{p{3cm}X X X}
\toprule
Domaine & Core & Smart & Optional / infrastructure \\
\midrule
Réception & UC-R01--UC-R07 & UC-R08 & UC-R09--UC-R12 \\
Stockage & UC-S01--UC-S05, UC-S07--UC-S08 & UC-S06, UC-S09--UC-S12 & UC-S13--UC-S14 \\
Commandes & UC-E01--UC-E09 & UC-E11, UC-E13, UC-E15 & UC-E10, UC-E12, UC-E14 \\
Yard & UC-Y01--UC-Y04 & UC-Y06--UC-Y07 & UC-Y05 \\
Ressources & UC-RES01--03, UC-RES05--06, UC-RES08 & UC-RES04 & UC-RES07, UC-RES09--10 \\
\bottomrule
\end{tabularx}
\end{table}

\chapter{Architecture fonctionnelle cible du Smart WMS}
\section{Principes de conception}
L'architecture cible est construite comme un système modulaire. Elle doit pouvoir s'appliquer à un entrepôt simple utilisant des terminaux RF, mais aussi à un site automatisé intégrant RFID, convoyeurs, AMR/AGV, WCS/WES ou BMS/EMS. Le modèle distingue donc les blocs fonctionnels stables, les services Smart activables par la donnée et les options dépendantes de l'infrastructure.

\section{Vue d'ensemble}
La figure~\ref{fig:architecture-fonctionnelle} redessine la vue fonctionnelle cible à partir de l'architecture validée dans le support de benchmark. Contrairement à une représentation trop compacte, elle sépare explicitement les systèmes externes, le cœur opérationnel, la gestion des tâches, les données de référence, le stockage des données, les ressources et infrastructures, l'intelligence décisionnelle et les appareils terrain.

\begin{landscape}
\begin{figure}[p]
\centering
\includegraphics[width=0.98\linewidth]{figures/chapitre4/architecture_fonctionnelle.pdf}
\caption{Architecture fonctionnelle cible du Smart WMS proposée.}
\label{fig:architecture-fonctionnelle}
\end{figure}
\end{landscape}

\section{Rôle des couches}
Les systèmes externes alimentent le WMS en commandes, statuts, confirmations, informations transport, production et référentiels. Le cœur opérationnel exécute les processus d'entrepôt : réception, stockage, inventaire, commandes, préparation, expédition et yard. La gestion des tâches transforme les décisions en actions exécutables et assure le suivi des statuts et exceptions.

Les données de référence ne sont pas traitées comme des use cases artificiels. Elles constituent une couche de support essentielle : produits, clients, fournisseurs, zones, règles, unités logistiques, utilisateurs et rôles. Le stockage est séparé en deux usages : l'OLTP pour l'exécution transactionnelle et l'OLAP/DWH pour l'historique, les KPI, l'analyse et les modèles prédictifs.

\section{Activation progressive}
L'architecture doit rester utile même lorsque toutes les technologies ne sont pas disponibles. La figure~\ref{fig:activation-progressive} formalise cette progressivité.

\begin{figure}[H]
\centering
\includegraphics[width=0.92\linewidth]{figures/chapitre4/activation_progressive.pdf}
\caption{Principe d'activation progressive des capacités Smart WMS.}
\label{fig:activation-progressive}
\end{figure}

\begin{table}[H]
\centering
\caption{Responsabilités des blocs fonctionnels de l'architecture cible}
\begin{tabularx}{\linewidth}{p{3.4cm}X X}
\toprule
Bloc & Responsabilité principale & Contribution au Smart WMS \\
\midrule
Cœur opérationnel & Exécuter réception, stock, commandes, yard et expédition & Stabiliser le socle métier indispensable \\
Gestion des tâches & Créer, affecter, suivre et traiter les exceptions & Transformer décisions et alertes en actions \\
Données de référence & Structurer produits, zones, règles, rôles et unités logistiques & Donner le contexte aux règles et optimisations \\
Ressources & Gérer main-d'œuvre, équipements, actifs et interfaces & Adapter le système au niveau réel d'automatisation \\
Intelligence & Prévoir, optimiser, détecter et recommander & Améliorer la qualité des décisions opérationnelles \\
Terrain & Capturer les données et guider l'exécution & Relier le numérique au flux physique \\
\bottomrule
\end{tabularx}
\end{table}

\chapter{Modélisation UML}
\section{Philosophie de modélisation}
La modélisation UML sert à rendre la cible fonctionnelle vérifiable. Les diagrammes ne doivent pas seulement illustrer le rapport ; ils doivent permettre de contrôler que chaque besoin métier possède un acteur, un processus, des interactions et, lorsque nécessaire, un composant applicatif.

\begin{figure}[H]
\centering
\includegraphics[width=0.92\linewidth]{figures/chapitre5/vue_modelisation_uml.pdf}
\caption{Vue de modélisation UML retenue pour le Smart WMS.}
\label{fig:vue-uml}
\end{figure}

\section{Diagrammes de cas d'utilisation}
Les diagrammes de cas d'utilisation sont organisés par domaine : vue globale, réception, stockage et inventaire, préparation et expédition, services transversaux, yard et ressources. La correction principale consiste à séparer visuellement les use cases Core, Smart et Optional. Les blocs Données de référence et Stockage des données ne sont pas représentés comme des cas d'utilisation déclenchés par un acteur ; ils relèvent du modèle de domaine et du diagramme de composants.

\section{Diagrammes d'activité}
Les diagrammes d'activité décrivent les flux de bout en bout : réception intelligente, stockage et inventaire, expédition, supervision énergétique et environnementale. Leur intérêt est de montrer les alternatives : règle métier standard lorsqu'aucune infrastructure avancée n'est disponible, puis activation OCR, caméra 3D, RFID, WCS/WES ou BMS/EMS lorsque le contexte le permet.

\section{Modèle de domaine et composants}
Le modèle de domaine doit couvrir les objets nécessaires à la cohérence métier : Warehouse, Zone, Location, Product, ProductRequirement, StockUnit, HandlingUnit, Lot, CustomerOrder, OrderLine, Task, Worker, Equipment, Device, Sensor, SensorReading, Alert et Recommendation. Le diagramme de composants traduira ensuite ces objets en services : services métier WMS, services Smart, intégrations externes, couche données et couche terrain.

Les sources PlantUML utilisées pour cette modélisation sont conservées en annexe afin de faciliter les itérations futures.

\chapter{Architecture applicative logique proposée}
\section{Objectif}
L'architecture applicative logique traduit la cible fonctionnelle en composants logiciels sans imposer une stack définitive. Elle sert à préparer une future implémentation tout en respectant le périmètre du stage : conception, modélisation et validation.

\begin{figure}[H]
\centering
\includegraphics[width=0.94\linewidth]{figures/chapitre6/architecture_applicative_logique.pdf}
\caption{Architecture applicative logique proposée pour le Smart WMS.}
\label{fig:architecture-applicative}
\end{figure}

\section{Positionnement du prototype SmartWMS AI Operations Agent}
Le dépôt local SmartWMS AI Operations Agent illustre une partie de cette architecture : une interface Next.js, une API FastAPI, un runner d'agent, des outils WMS en lecture seule, un moteur de décisions déterministes et un retriever de documents synthétiques. Cette maquette ne constitue pas un WMS complet. Elle démontre surtout comment une couche d'aide à la décision peut consommer des faits structurés sans inventer la logique métier.

\section{Principes de sûreté}
Les règles de décision doivent rester testables hors du modèle de langage. Dans la maquette, les calculs de disponibilité, position de stock, demande ouverte, risque de rupture, priorité de réapprovisionnement et quantité de réapprovisionnement sont centralisés dans un service métier déterministe. Le modèle de langage explique les résultats ; il ne crée pas de mouvements de stock et ne prend pas de décision d'achat autonome.

\chapter{Validation, limites et perspectives}
\section{Méthode de validation}
La validation proposée repose sur des scénarios métier plutôt que sur un déploiement industriel. Chaque scénario doit être relié à un problème, à un ensemble de use cases, à un diagramme UML et à un composant applicatif.

\begin{figure}[H]
\centering
\includegraphics[width=0.88\linewidth]{figures/chapitre7/chaine_validation.pdf}
\caption{Chaîne de validation proposée pour l'architecture Smart WMS.}
\label{fig:validation}
\end{figure}

\section{Limites}
Le benchmark reste limité aux informations disponibles et ne remplace pas une étude détaillée de licences, coûts, performances ou contraintes réglementaires. La modélisation UML n'a pas encore été validée sur un entrepôt réel. Les capacités IA restent conceptuelles ou démontrées par une maquette contrôlée, avec données synthétiques.

\section{Perspectives}
Les perspectives naturelles sont la simulation de scénarios réels, la collecte de jeux de données entrepôt, la création d'un dashboard de supervision, l'intégration d'un IoT gateway, l'évaluation de modèles de prévision, puis la connexion progressive avec ERP, WCS/WES ou BMS/EMS. Une attention particulière devra être portée à la qualité des données, à la sécurité, à la supervision humaine et à la traçabilité des recommandations.

\chapter*{Conclusion générale}
Ce stage a permis de structurer une vision cohérente de la transformation d'un WMS traditionnel vers un Smart WMS. Le travail ne s'est pas limité à identifier des technologies avancées ; il a consisté à replacer chaque capacité dans un besoin métier, une condition d'activation et une architecture fonctionnelle claire.

Le benchmark a montré trois principes complémentaires : l'unification de l'exécution, la richesse du modèle métier et l'orchestration des tâches. L'analyse des besoins a ensuite permis de formaliser les use cases Core, Smart et Optional. L'architecture fonctionnelle cible positionne les systèmes externes, le cœur opérationnel, les tâches, les données, les ressources, l'intelligence et les terminaux terrain dans une structure modulaire.

La suite du travail consiste à consolider la modélisation UML, valider les scénarios avec des acteurs métier et transformer progressivement l'architecture logique en prototype plus complet. L'apport principal du stage réside dans cette clarification : un Smart WMS robuste n'est pas un WMS auquel on ajoute de l'IA partout, mais un système capable d'utiliser les bonnes données, au bon endroit, pour améliorer les décisions et l'exécution réelle.


\appendix
\chapter{Sources PlantUML}
Les fichiers PlantUML du projet sont conservés dans \texttt{annexes/plantuml}. Ils servent de base aux diagrammes de cas d'utilisation et d'activité et doivent être rendus avec PlantUML lors d'une itération outillée.

\bibliographystyle{plain}
\bibliography{bibliography/references}
\end{document}
